\subsection{Parameters and Exceptional Values}

When a matrix contains a parameter, the equation
\[
\det A(t)=0
\]
identifies values at which a relation among rows or columns appears.

\begin{examplebox}[A quadratic parameter condition]
Let
\[
A(t)=\begin{pmatrix}t&1\\4&t\end{pmatrix}.
\]
Then
\[
\det A(t)=t^2-4=(t-2)(t+2).
\]
Thus the determinant vanishes for $t=2$ or $t=-2$. At $t=2$ the second row is twice the first; at $t=-2$ it is $-2$ times the first.
\end{examplebox}

\begin{examplebox}[A row relation in three dimensions]
For
\[
B(s)=\begin{pmatrix}
1&0&1\\0&1&1\\1&s&2
\end{pmatrix},
\]
we obtain
\[
\det B(s)=1-s.
\]
At $s=1$ the third row satisfies
\[
(1,1,2)=(1,0,1)+(0,1,1).
\]
\end{examplebox}

\begin{summarybox}[How to read a zero determinant]
Solve $\det A=0$, then inspect the exceptional matrices for an explicit relation among rows or columns. The roots are not merely numerical answers; they locate the loss of independence.
\end{summarybox}
